\subsection{Direct Proof}
\label{subsections:techniques:universal:direct}



\paragraph{Description.}
A direct proof proves a mathematical statement by reasoning straight from %
what is assumed to what is claimed. %
Precisely, it proves a universal statement, $\forall s \in S,\ P(s)$, %
by letting $s \in S$ be arbitrary and deducing that $P(s)$ is true. %
Because $s$ is arbitrary, no specific value is used, %
so the argument holds for every element of $S$. %
Most subsequent proofs in this document are common applications or %
variations of a direct proof. %
Generally, a direct proof mirrors the following format:
\begin{enumerate}{
    \item \textbf{Setup.} %
        Let $s \in S$ be arbitrary, assuming any given hypotheses about $s$.
    \item \textbf{Verification.} %
        Show that $P(s)$ is true.
    \item \textbf{Conclusion.} %
        Conclude the proof, optionally by summarizing what has been found.
}
\end{enumerate}



\paragraph{Advice.}
\begin{itemize}{
    \item Use this proof as the default approach for a universal statement, %
        $\forall s \in S,\ P(s)$.
    \item Do not assume a specific value for $s$.
    \item A direct proof may evolve into another proof as it is written, %
        for instance when a step requires considering cases.
}
\end{itemize}



\paragraph{Example.}
\begin{exercise}{
    Prove that if $z \in \Z$ is even, %
    then $\exists m \in \Z$ such that $z^2 = 4m$.
}
\end{exercise}
\begin{proof}{
    This statement can be proven directly.

    \begin{annotatedpanel}[Setup.]{
        Let $z \in \Z$ such that $z$ is even. %
        Then $z = 2k$ with $k \in \Z$.
    }{
        \textbf{Step 1.} Let $s \in S$ be arbitrary, assuming any given %
        hypotheses. %
        Here, $S = \Z$, $s = z$, and the hypothesis is that $z$ is even.
    }
    \end{annotatedpanel}

    \begin{annotatedpanel}[Verification.]{
        Then $z^2 = (2k)^2 = 4k^2$. %
        Let $m \in \Z$ such that $m = k^2$. %
        Then $z^2 = 4k^2 = 4m$.
    }{
        \textbf{Step 2.} Show that $P(s)$ is true. %
        Here, $P(z)$ is the statement $\exists m \in \Z$ such that %
        $z^2 = 4m$, and $m = k^2$ satisfies it.
    }
    \end{annotatedpanel}

    \begin{annotatedpanel}[Conclusion.]{
        Thus, $\exists m \in \Z$ such that $z^2 = 4m$. %
        Therefore, the statement holds. \qedhere
    }{
        \textbf{Step 3.} Conclude the proof. %
        Here, the summary states what has been found: %
        $z^2$ is a multiple of $4$.
    }
    \end{annotatedpanel}
}
\end{proof}